\section{Реализация алгоритмов}

Реализация алгоритмов выполнена на языке программирования C\texttt{++} с использованием стандартной библиотеки \cite{cppreference}. Выбор языка обусловлен его высокой производительностью и отсутствием накладных расходов, что важно для точного измерения времени выполнения.

\subsection{Пузырьковая сортировка}

В листинге~\ref{code:bubble} представлена реализация пузырьковой сортировки с оптимизацией. Переменная \texttt{swapped} используется для досрочного завершения алгоритма: если за весь проход не было ни одного обмена, массив уже отсортирован. Параметром функции является вектор \texttt{arr}, который передаётся по ссылке.

\begin{listing}[h]
\caption{Пузырьковая сортировка на C++}
\label{code:bubble}
\begin{minted}[frame=single,linenos]{cpp}
void bubbleSort(vector<int>& arr) {
    int n = arr.size();
    bool swapped;
    for (int i = 0; i < n-1; i++) {
        swapped = false;
        for (int j = 0; j < n-i-1; j++) {
            if (arr[j] > arr[j+1]) {
                swap(arr[j], arr[j+1]);
                swapped = true;
            }
        }
        if (!swapped) break;
    }
}
\end{minted}
\end{listing}

\subsection{Быстрая сортировка}

Реализация быстрой сортировки состоит из двух функций. Функция \texttt{partition} (листинг~\ref{code:partition}) выбирает опорный элемент \texttt{pivot} и разделяет массив на две части. Параметры \texttt{low} и \texttt{high} задают границы сортировки. Опорный элемент \texttt{pivot} выбирается как последний элемент массива (\texttt{arr[high]}). Переменная \texttt{i} используется для отслеживания границы меньших элементов.

\begin{listing}[h]
\caption{Функция разделения (partition)}
\label{code:partition}
\begin{minted}[frame=single,linenos]{cpp}
int partition(vector<int>& arr, int low, int high) {
    int pivot = arr[high];
    int i = low - 1;
    for (int j = low; j < high; j++) {
        if (arr[j] <= pivot) {
            i++;
            swap(arr[i], arr[j]);
        }
    }
    swap(arr[i+1], arr[high]);
    return i+1;
}
\end{minted}
\end{listing}

Функция \texttt{quickSort} (листинг~\ref{code:quicksort}) рекурсивно вызывает себя для левой и правой частей. Параметр \texttt{pi} — это индекс опорного элемента \texttt{pivot}, возвращаемый функцией \texttt{partition}.

\begin{listing}[h]
\caption{Быстрая сортировка (quickSort)}
\label{code:quicksort}
\begin{minted}[frame=single,linenos]{cpp}
void quickSort(vector<int>& arr, int low, int high) {
    if (low < high) {
        int pi = partition(arr, low, high);
        quickSort(arr, low, pi-1);
        quickSort(arr, pi+1, high);
    }
}
\end{minted}
\end{listing}

В приведённой реализации опорным элементом \texttt{pivot} выбирается последний элемент массива. Это простой, но не оптимальный выбор. Для улучшения производительности можно использовать выбор среднего элемента или случайного опорного элемента.